% VERSION VERIFIEE -- 19 juin 2026
% !TEX root = main.tex
% Annexe « Enseignements » -- version finale
% À inclure dans la partie annexes, après l'appel à \appendix et avant
% \SortirDesAnnexes, par exemple : \include{annexe_enseignement_final}

\chapter{Enseignements}
\label{ann:enseignement}

\begingroup

% -----------------------------------------------------------------------------
% Charte locale : nuances de gris, cohérentes avec les autres annexes.
% -----------------------------------------------------------------------------
\colorlet{EnsBlack}{black!88}
\colorlet{EnsDark}{black!70}
\colorlet{EnsMid}{black!42}
\colorlet{EnsLight}{black!15}
\colorlet{EnsPale}{black!5}

% Liens cliquables, sans faire apparaître les URL dans le corps du texte.
\hypersetup{colorlinks=true,allcolors=EnsBlack,pdfborder={0 0 0}}

\newcommand{\EnsBadge}[1]{%
  \begingroup
  \setlength{\fboxsep}{2.5pt}%
  \colorbox{EnsPale}{%
    \textcolor{EnsDark}{\scriptsize\bfseries\textsc{#1}}%
  }%
  \endgroup
}

\newcommand{\EnsAnnee}[4]{%
  \multicolumn{2}{@{}l@{}}{%
    \parbox{\textwidth}{%
      \vspace{.35\baselineskip}%
      \noindent
      {\color{EnsBlack}\large\bfseries Année universitaire #1}%
      \hfill
      {\scriptsize\color{EnsDark}%
        \href{#2}{\textsc{Radar} Ayana #3}%
      }%
      \par\vspace{.16\baselineskip}%
      {\footnotesize\color{EnsDark}\textbf{Volume personnel total :} #4}%
      \par\vspace{.22\baselineskip}%
      {\color{EnsLight}\hrule height .7pt}%
    }%
  }\\*[.85\baselineskip]
}

\newcommand{\EnsEntree}[4]{%
  \begin{minipage}[t]{\linewidth}
    \raggedleft
    #1
  \end{minipage}%
  &
  \begin{minipage}[t]{\linewidth}
    \raggedright
    #2\par\vspace{.42\baselineskip}%
    #3\par\vspace{.42\baselineskip}%
    #4
  \end{minipage}%
  \\[1.35\baselineskip]
}

% \begin{center}
%   \setlength{\fboxsep}{9pt}%
%   \colorbox{EnsPale}{%
%     \parbox{.91\textwidth}{%
%       \small
%       Cette annexe recense les enseignements effectivement assurés par Louis
%       Hauseux à
%       l'\href{https://univ-cotedazur.fr/}{Université Côte d'Azur}, dans le
%       \href{https://univ-cotedazur.eu/msc/msc-data-science-and-artificial-intelligence}{%
%       \emph{MSc Data Science and Artificial Intelligence}}. Les rapports
%       d'activité Inria étant structurés par année civile, les interventions ont
%       ici été replacées dans leur année universitaire et présentées par ordre
%       antéchronologique. Pour les unités d'enseignement coassurées, le
%       \emph{volume personnel} désigne uniquement les séances données par Louis
%       Hauseux; il est distingué du \emph{volume complet de l'UE}. Les cours
%       magistraux (CM), travaux dirigés ou pratiques (TD) et volumes de
%       coenseignement sont donc distingués explicitement, afin de ne pas
%       attribuer à un seul intervenant le volume collectif d'une unité
%       d'enseignement.
%     }%
%   }
% \end{center}

\vspace{.45\baselineskip}
\begin{center}
  \small
  \EnsBadge{CM}\enspace cours magistral
  \qquad
  \EnsBadge{TD}\enspace travaux dirigés ou pratiques
  \qquad
  \EnsBadge{hETD}\enspace heure équivalent TD
\end{center}

\begin{center}
  \footnotesize\color{EnsDark}
  Convention employée :
  \(1\,\mathrm{h\,CM}=1{,}5\,\mathrm{hETD}\) et
  \(1\,\mathrm{h\,TD}=1\,\mathrm{hETD}\).
\end{center}

\vspace{.65\baselineskip}
\begin{center}
  {\color{EnsBlack}\bfseries Synthèse des heures effectivement assurées}\par
  \vspace{.45\baselineskip}
  \renewcommand{\arraystretch}{1.28}%
  \setlength{\tabcolsep}{9pt}%
  \begin{tabular}{@{}lrrrr@{}}
    \toprule
    \rowcolor{EnsPale}
    \textbf{Année universitaire}
      & \textbf{CM}
      & \textbf{TD}
      & \textbf{Présence}
      & \textbf{hETD} \\
    \midrule
    2025--2026 & 8 h 15  & 50 h 15 & 58 h 30  & 62{,}625 \\
    2024--2025 & 10 h 30 & 40 h 30 & 51 h     & 56{,}25  \\
    2023--2024 & 10 h 30 & 19 h 30 & 30 h     & 35{,}25  \\
    \midrule
    \textbf{Total} & \textbf{29 h 15} & \textbf{110 h 15}
      & \textbf{139 h 30} & \textbf{154{,}125} \\
    \bottomrule
  \end{tabular}
  \par\vspace{.45\baselineskip}
  \parbox{.88\textwidth}{%
    \footnotesize\color{EnsDark}
    % Les volumes ci-dessus correspondent aux heures d'enseignement en présence
    % effectivement assurées par Louis Hauseux. Ils n'incluent ni la préparation
    % des supports et des évaluations, ni les corrections, ni le suivi
    % pédagogique hors séance.
  }
\end{center}

\clearpage

\vspace{.45\baselineskip}
\small
\setlength{\LTpre}{0pt}
\setlength{\LTpost}{0pt}
\setlength{\tabcolsep}{0pt}
\renewcommand{\arraystretch}{1.04}

\begin{longtable}{%
  @{}%
  >{\raggedleft\arraybackslash\color{EnsDark}}p{.175\textwidth}%
  @{\hspace{1.00em}}%
  !{\color{EnsMid}\vrule width .55pt}%
  @{\hspace{1.20em}}%
  >{\raggedright\arraybackslash\color{black}}p{.735\textwidth}%
  @{}%
}

\endfirsthead
\multicolumn{2}{@{}r@{}}{%
  \scriptsize\color{EnsMid}\emph{Enseignements --- suite}%
}\\[.65\baselineskip]
\endhead

% =============================================================================
% 2025--2026
% =============================================================================
\EnsAnnee{2025--2026}{https://radar.inria.fr/report/2025/ayana/index.html}{2025}{8 h 15 de CM + 50 h 15 de TD = 58 h 30 devant les étudiants (62{,}625 hETD)}

\EnsEntree{%
  \textbf{Févr.--avr. 2026}\par\vspace{.35\baselineskip}
  \EnsBadge{M2}\enspace\EnsBadge{TD}\par\vspace{.35\baselineskip}
  2 ECTS 
}{%
  {\large\bfseries
  \href{https://univ-cotedazur.eu/msc/msc-data-science-and-artificial-intelligence/program/second-year}{%
  \emph{Statistical Learning Theory}}}%
}{%
  \textbf{Formation et établissement.} Deuxième année du
  \emph{MSc Data Science and Artificial Intelligence}, parcours
  \emph{Core AI}, Université Côte d'Azur. Le cours est sous la
  responsabilité de
  \href{https://www-sop.inria.fr/members/Mathieu.Carriere/}{Mathieu \nom{Carrière}}. %,  chargé de recherche Inria au sein de l'équipe DataShape et titulaire d'une chaire 3IA Côte d'Azur. 
  % La maquette présente cet enseignement comme une
  % introduction aux principes de la théorie statistique de l'apprentissage,
  % associant développements théoriques et séances pratiques de programmation.
}{%
  \textbf{Part personnelle.} \textbf{Huit séances (sur huit) de 1 h 30 de TD},
  réparties de février à avril 2026, soit \textbf{12 heures de TD}. 
  % Les séances assurées par Louis Hauseux
  % accompagnent le cours de Mathieu \nom{Carrière} par des exercices et mises en
  % œuvre numériques consacrés à l'analyse des régimes de fonctionnement, des
  % capacités et des garanties des modèles d'apprentissage statistique.
  % \par\smallskip
  % \textbf{Volume complet de l'UE.} La
  % \href{https://univ-cotedazur.eu/msc/msc-data-science-and-artificial-intelligence/program/second-year}{maquette publiée}
  % indique \textbf{24 heures et 2 ECTS}. Elle précise que l'UE combine cours et
  % séances pratiques, sans publier leur ventilation horaire détaillée. Le volume
  % retenu dans cette annexe correspond donc exclusivement aux \textbf{12 heures
  % de TD effectivement assurées par Louis Hauseux}.
}

\EnsEntree{%
  \textbf{Sept.--nov. 2025}\par\vspace{.35\baselineskip}
  \EnsBadge{M2}\enspace\EnsBadge{CM + TD}\par\vspace{.35\baselineskip}
  2 ECTS
}{%
  {\large\bfseries
  \href{https://univ-cotedazur.eu/msc/msc-data-science-and-artificial-intelligence/program/second-year}{%
  \emph{Statistical Analysis of Graphs (Networks)}}}%
}{%
  \textbf{Formation et établissement.} Deuxième année du
  \emph{MSc Data Science and Artificial Intelligence}, parcours
  \emph{Core AI}, Université Côte d'Azur. 
  % L'intitulé abrégé employé dans les
  % rapports Inria est \emph{Statistical Analysis of Networks}. 
  Co-enseignement avec Konstantin \nom{Avrachenkov} et Rémy \nom{Sun}.
}{%
  \textbf{Part personnelle.} \textbf{Cinq séances et demie sur huit}, soit
  \textbf{8 h 15 de CM et 8 h 15 de TD} : 16 h 30 devant les étudiants,
  correspondant à \textbf{20{,}625 hETD}. 
  % La répartition administrative
  % attribue cinq séances et demie à Louis Hauseux, une séance et demie à
  % Konstantin \nom{Avrachenkov} et une séance à Rémy \nom{Sun}, consacrée à
  % une introduction aux réseaux de neurones sur graphes.
  % \par\smallskip
  % \textbf{Volume complet de l'UE.} Huit créneaux de trois heures, chacun
  % composé administrativement de 1 h 30 de CM et 1 h 30 de TD, soit
  % \textbf{12 h de CM + 12 h de TD = 24 h en présence = 30 hETD}.
  % L'intervention personnelle porte notamment sur les rappels d'algèbre
  % linéaire et de probabilités, les modèles de graphes aléatoires, les graphes
  % géométriques et le Single-Linkage, les méthodes spectrales, les méthodes de
  % coupe, la modularité et leur mise en œuvre dans des notebooks Python.
}

\EnsEntree{%
  \textbf{Oct.--déc. 2025}\par\vspace{.35\baselineskip}
  \EnsBadge{M1}\enspace\EnsBadge{TD}\par\vspace{.35\baselineskip}
  6 ECTS 
}{%
  {\large\bfseries
  \href{https://univ-cotedazur.eu/msc/msc-data-science-and-artificial-intelligence/program/first-year}{%
  \emph{Statistical Inference Theory}}}%
}{%
  \textbf{Formation et établissement.} Première année du
  \emph{MSc Data Science and Artificial Intelligence}, Université Côte
  d'Azur. % L'intitulé abrégé des rapports Inria est \emph{Statistical Inference}. 
  Le cours magistral est assuré par Vincent \nom{Vandewalle}.
}{%
  \textbf{Part personnelle.} \textbf{30 heures de TD}. 
  % Louis Hauseux a préparé et animé les séances
  % d'exercices et de programmation, produit les énoncés et corrigés associés et
  % participé à la préparation, à la surveillance et à la correction des
  % évaluations. Les TD portent notamment sur l'estimation, les lois empiriques,
  % le bootstrap, les tests statistiques et les modèles de régression.
  % \par\smallskip
  % \textbf{Volume complet de l'UE.} La
  % \href{https://univ-cotedazur.eu/msc/msc-data-science-and-artificial-intelligence/program/first-year}{maquette publiée}
  % indique 60 heures et 6 ECTS; les 30 heures ci-dessus correspondent
  % exclusivement aux TD assurés par Louis Hauseux.
}

% =============================================================================
% 2024--2025
% =============================================================================
\EnsAnnee{2024--2025}{https://radar.inria.fr/report/2024/ayana/index.html}{2024}{10 h 30 de CM + 40 h 30 de TD = 51 h devant les étudiants (56{,}25 hETD)}

\EnsEntree{%
  \textbf{Oct.--déc. 2024}\par\vspace{.35\baselineskip}
  \EnsBadge{M2}\enspace\EnsBadge{CM + TD}\par\vspace{.35\baselineskip}
  2 ECTS
}{%
  {\large\bfseries
  \href{https://univ-cotedazur.eu/msc/msc-data-science-and-artificial-intelligence/program/second-year}{%
  \emph{Statistical Analysis of Graphs (Networks)}}}%
}{%
  \textbf{Formation et établissement.} Deuxième année du
  \emph{MSc Data Science and Artificial Intelligence}, parcours
  \emph{Core AI}, Université Côte d'Azur. Co-enseignement avec Konstantin
  \nom{Avrachenkov}. %; les rapports Inria emploient l'intitulé abrégé \emph{Statistical Analysis of Networks}.
}{%
  \textbf{Part personnelle.} \textbf{Sept séances sur huit}, soit
  \textbf{10 h 30 de CM et 10 h 30 de TD} : 21 heures devant les étudiants,
  correspondant à \textbf{26{,}25 hETD}. 
  % Une séance, consacrée aux
  % méthodes d'échantillonnage par chaînes de Markov, a été assurée par
  % Konstantin \nom{Avrachenkov}.
  % \par\smallskip
  % \textbf{Volume complet de l'UE.} \textbf{12 h de CM + 12 h de TD}, soit
  % 24 heures en présence et 30 hETD. Les séances assurées par Louis Hauseux
  % couvrent les fondements de la théorie des graphes, les modèles
  % d'Erdős--Rényi, à blocs stochastiques et géométriques, la percolation, les
  % indices de centralité, PageRank, les méthodes spectrales, la modularité et la
  % détection de communautés, avec des travaux pratiques en Python.
}

\EnsEntree{%
  \textbf{Oct.--déc. 2024}\par\vspace{.35\baselineskip}
  \EnsBadge{M1}\enspace\EnsBadge{TD}\par\vspace{.35\baselineskip}
  6 ECTS 
}{%
  {\large\bfseries
  \href{https://univ-cotedazur.eu/msc/msc-data-science-and-artificial-intelligence/program/first-year}{%
  \emph{Statistical Inference Theory}}}%
}{%
  \textbf{Formation et établissement.} Première année du
  \emph{MSc Data Science and Artificial Intelligence}, Université Côte
  d'Azur. Le cours magistral est assuré par Vincent \nom{Vandewalle}.
  %; les rapports Inria utilisent l'intitulé abrégé \emph{Statistical Inference}.
}{%
  \textbf{Part personnelle.} \textbf{30 heures de TD}. %, sans cours magistral, soit \textbf{30 hETD}.
  % , conformément au volume déclaré dans le rapport
  % d'activité Ayana 2024. Louis Hauseux a assuré les séances d'exercices et de
  % programmation, préparé des supports et corrigés, suivi les travaux à rendre
  % et contribué aux examens. Les applications mobilisent notamment R et Python.
  % \par\smallskip
  % \textbf{Volume complet de l'UE.} La maquette publiée indique 60 heures et
  % 6 ECTS; le volume de 30 heures reporté ici est uniquement le service de TD de
  % Louis Hauseux, et non le volume intégral de l'UE.
}

% =============================================================================
% 2023--2024
% =============================================================================
\EnsAnnee{2023--2024}{https://radar.inria.fr/report/2023/ayana/index.html}{2023}{10 h 30 de CM + 19 h 30 de TD = 30 h devant les étudiants (35{,}25 hETD)}

\EnsEntree{%
  \textbf{Sept.--déc. 2023}\par\vspace{.35\baselineskip}
  \EnsBadge{M2}\enspace\EnsBadge{CM + TD}\par\vspace{.35\baselineskip}
  2 ECTS
}{%
  {\large\bfseries
  \href{https://univ-cotedazur.eu/msc/msc-data-science-and-artificial-intelligence/program/second-year}{%
  \emph{Statistical Analysis of Graphs (Networks)}}}%
}{%
  \textbf{Formation et établissement.} Deuxième année du
  \emph{MSc Data Science and Artificial Intelligence}, parcours
  \emph{Core AI}, Université Côte d'Azur. Première édition assurée pendant la
  thèse, en co-enseignement avec Konstantin \nom{Avrachenkov}. %; l'intitulé du  rapport d'activité est \emph{Statistical Analysis of Networks}.
}{%
  \textbf{Part personnelle.} \textbf{Sept séances sur huit}, soit
  \textbf{10 h 30 de CM et 10 h 30 de TD} : 21 heures devant les étudiants,
  correspondant à \textbf{26{,}25 hETD}. 
  %Konstantin \nom{Avrachenkov} a assuré la dernière séance, consacrée aux méthodes d'échantillonnage.
  % Outre
  % l'enseignement, Louis Hauseux a préparé les supports de cours, les notebooks
  % pratiques, les devoirs et l'évaluation finale.
}

\EnsEntree{%
  \textbf{Mars 2024}\par\vspace{.35\baselineskip}
  \EnsBadge{M1}\enspace\EnsBadge{TD}\par\vspace{.35\baselineskip}
  Intervention ponctuelle
}{%
  {\large\bfseries
  \href{https://univ-cotedazur.eu/msc/msc-data-science-and-artificial-intelligence/program/first-year}{%
  \emph{Statistical Inference Theory}}}%
}{%
  \textbf{Formation et établissement.} Première année du
  \emph{MSc Data Science and Artificial Intelligence}, Université Côte
  d'Azur. Intervention pour des séances de révision intégrale pour les dix élèves n'ayant pas validé la première session du cours de Vincent
  \nom{Vandewalle}.
}{%
  \textbf{Part personnelle.} \textbf{Six séances de 1 h 30}, soit
  \textbf{9 heures de TD}, (comptabilisés \textbf{9 hETD}). Le
  \href{https://overleaf.irisa.fr/read/wjmztthswkwy}{support préparé pour ces
  séances} rassemble rappels de cours, exercices et éléments de correction.
  % \par\smallskip
  % Cette intervention est volontairement présentée comme ponctuelle : les neuf
  % heures correspondent exactement aux séances données par Louis Hauseux en
  % mars 2024, et non à l'ensemble des TD ni aux 60 heures de l'UE complète.
}

\end{longtable}

% \vspace{.15\baselineskip}
% \enlargethispage{3\baselineskip}
% \noindent
% {\scriptsize\color{EnsDark}
%   \textbf{Sources et méthode.}
%   Les informations ont été recoupées dans les
%   \href{https://radar.inria.fr/report/2023/ayana/index.html}{rapports
%   \textsc{Radar} Ayana 2023},
%   \href{https://radar.inria.fr/report/2024/ayana/index.html}{2024} et
%   \href{https://radar.inria.fr/report/2025/ayana/index.html}{2025}, les
%   \href{https://univ-cotedazur.eu/msc/msc-data-science-and-artificial-intelligence/program/first-year}{maquettes de M1} et
%   \href{https://univ-cotedazur.eu/msc/msc-data-science-and-artificial-intelligence/program/second-year}{de M2},
%   la \href{https://www-sop.inria.fr/members/Mathieu.Carriere/}{notice institutionnelle de Mathieu \nom{Carrière}},
%   ainsi que les relevés d'intervention de l'auteur. Les rapports Inria étant
%   annuels civils, les enseignements sont ici reclassés par année universitaire.
% \par}

\endgroup
